\section*{Contribution and Novelty}
\label{sec:contribution-novelty}
Within the smoothly regularized three-field remodeling framework (density $\rho$, stimulus $S$, and log-fabric $\mathbf{L}$), we contribute two robustness mechanisms that enable stable FE simulations in stiff coupling regimes. First, we introduce a numerically robust fabric-guided orthotropy pipeline based on stabilized Sylvester spectral projectors, smoothly blended to the isotropic and transversely isotropic limits. This enforces the partition of unity and avoids basis flips near eigenvalue degeneracy, keeping the constitutive evaluation and the trace-free fabric target well defined as anisotropy evolves.
Second, we stabilize the fully implicit coupled update by a partitioned block Gauss--Seidel strategy with safeguarded Anderson (Pulay/DIIS) acceleration, including adaptive restarts and step limiting, which delivers robust convergence where plain Picard iteration can stall. \rev{Beyond these numerical contributions, we introduce a modelling element: a compartment-aware reference stimulus $\hat{\psi}_{\mathrm{ref}}(\widetilde{\rho})$ that blends distinct trabecular and cortical mechanostat set points via a smooth density-dependent transition. This reflects accumulating experimental evidence that remodeling thresholds are not universal constants but differ between bone compartments, and we show that a non-unity cortical--trabecular ratio is necessary---within our framework and loading assumptions---to recover anatomically plausible density distributions in the proximal femur case study.} Finally, we demonstrate the framework on benchmark tests and a population-CT-derived proximal femur study.
